\section*{Chapter \ref{ch:iv1:counting-and-symmetry} --- Counting and Symmetry}
\markboth{\protect\omTocMark{}{Answers: Counting and Symmetry}}{\protect\omTocMark{}{Answers: Counting and Symmetry}}

\begin{solution}{iq:iv1:counting-and-symmetry:1}
There are $7! = 5040$ seatings, and $3600$ keep the two apart.

Fix one person's chair to remove the rotations, and then seat the other seven in order. For the second count, subtract the seatings with the two together. Glue them into a block: that leaves seven objects around the table, which gives $6!$ arrangements, and the block has 2 internal orders. So $5040 - 2 \cdot 720 = 3600$.

\iqfollow{What fraction is that?}{$5/7$. Seen from one of the two, the other sits in one of seven chairs at random, and two of those are adjacent.}
\end{solution}

\begin{solution}{iq:iv1:counting-and-symmetry:2}
There are $34\,650$ arrangements, of which $7350$ have no two S's together.

There are 11 letters, with four I's, four S's and two P's, so the multinomial coefficient\index{multinomial coefficient} gives $11!/(4!\,4!\,2!) = 34\,650$.

For the second count, place the S's last. The seven other letters can be arranged in $7!/(4!\,2!) = 105$ ways. They leave 8 gaps, counting both ends, and the four S's go into four different gaps, in $\binom{8}{4} = 70$ ways. The product is $7350$, or $7/33$ of the total.

\iqwrong{Subtracting ``arrangements with an SS block'' from the total. Blocks overlap in many ways, and the gaps method avoids all of them.}
\end{solution}

\begin{solution}{iq:iv1:counting-and-symmetry:3}
There are 252 paths, and 132 avoid the point.

A path is a string of five rights and five ups, so there are $\binom{10}{5} = 252$. Those through $(2, 2)$ split into two independent parts, $\binom{4}{2} = 6$ ways to reach it and $\binom{6}{3} = 20$ ways to go on. So 120 pass through it, and $252 - 120 = 132$ avoid it.

\iqfollow{How would you code it for many blocked points?}{Fill a table by dynamic programming: each cell holds the sum of the cells to its left and below, and zero if blocked.}
\end{solution}

\begin{solution}{iq:iv1:counting-and-symmetry:4}
There are 286 ways with empty boxes allowed, and 84 ways with none empty.

Use stars and bars\index{stars and bars}. Line up the 10 coins and put 3 dividers among them. That is a choice of 3 places out of 13, so $\binom{13}{3} = 286$. If no box may be empty, put one coin in each first. That leaves 6 coins to share in any way: $\binom{9}{3} = 84$.

\iqwrong{$4^{10}$, which is the answer for labelled coins. Identical coins make the outcome depend only on how many land in each box.}
\end{solution}

\begin{solution}{iq:iv1:counting-and-symmetry:5}
There are 315 committees. The shortcut gives 600, because it counts a committee once for each woman in it.

Count the complement\index{complement counting}: of the $\binom{11}{4} = 330$ committees, $\binom{6}{4} = 15$ are all men, and $330 - 15 = 315$.

The wrong method names one woman as ``the'' woman and then chooses the remaining three, giving $5 \binom{10}{3} = 600$. A committee with $w$ women is produced $w$ times, once for each of them. Adding up the number of women over all committees gives exactly 600.

\iqwhy{Whenever a count looks too big, ask how many times each object is produced. That question finds most counting errors.}
\end{solution}

\begin{solution}{iq:iv1:counting-and-symmetry:6}
There are 252 different results, and the probability of a full house is $300/7776 \approx 0.0386$.

A result ignoring order is a multiset of 5 values from 6, so stars and bars\index{stars and bars} gives $\binom{10}{5} = 252$.

The results are not equally likely, so you cannot divide by 252. A full house, say 3, 3, 3, 5, 5, arises from $\binom{5}{2} = 10$ ordered rolls, while 1, 2, 3, 4, 5 arises from 120 and five sixes from just one. So count ordered rolls: 6 choices for the triple, 5 for the pair, 10 orders, which is 300 out of $6^5 = 7776$.

\iqwrong{Dividing by 252 to get $30/252 \approx 0.119$. Treating unordered results as equally likely is the most common error in this chapter's topic.}
\end{solution}

\begin{solution}{iq:iv1:counting-and-symmetry:7}
Both are $3/8$.

Use symmetry\index{exchangeability}. Every ordering of the eight balls is equally likely, so every position has the same chance of holding a given ball. Each of the 8 balls is equally likely to be last, and 3 of them are red. The same holds for the second position, and for any other.

\iqwrong{Conditioning through all the earlier draws. This works, but it takes a page, and the symmetry makes the order in which balls come out irrelevant. The first ball you see does change the conditional chance for the second, but before you see anything, each position has the same chance.}
\end{solution}

\begin{solution}{iq:iv1:counting-and-symmetry:8}
There are 8 with three painted faces, 96 with two, 384 with one and 512 with none. On average a unit cube has $0.6$ painted faces.

The corners give the 8. Each of the 12 edges has $10 - 2 = 8$ inner cubes, which gives 96. Each of the 6 faces has $8^2 = 64$ interior squares, which gives 384. The inside is a block of $8^3 = 512$. These add to 1,000, which is the check.

For the average, count painted faces from the outside: the big cube has 6 faces of 100 unit squares each, which makes 600 faces. Spread over 1,000 cubes that is $0.6$.

\iqfollow{And for a side of $n$?}{$6n^2/n^3 = 6/n$ faces on average.}
\end{solution}

\begin{solution}{iq:iv1:counting-and-symmetry:9}
Exactly 3 has probability $\approx 0.0177$, or about 1 in 57, and at least 3 has probability $\approx 0.0186$, or about 1 in 54.

Count the draws by how they split between your 6 numbers and the other 43: the hypergeometric distribution\index{hypergeometric distribution}. There are $\binom{49}{6} = 13\,983\,816$ draws. Exactly 3 matches come from $\binom{6}{3}\binom{43}{3} = 20 \cdot 12\,341 = 246\,820$ draws. Summing over 3, 4, 5 and 6 matches gives $0.0186$.

\iqwrong{Using a binomial with $p = 6/49$. Draws are made without replacement, so the matches are not independent. The binomial gives a value near $0.025$, too high.}
\end{solution}

\enlargethispage{-5\baselineskip}%
\begin{solution}{iq:iv1:counting-and-symmetry:10}
There are 266.

Use inclusion-exclusion\index{inclusion-exclusion}. Count the multiples of each prime and subtract, then add back those subtracted twice, then remove those added back too often:
\[
1000 - (500 + 333 + 200) + (166 + 100 + 66) - 33 = 266.
\]
The middle terms are multiples of 6, 10 and 15, and the last term is multiples of 30.

A quick check is by density. A share $\tfrac12 \cdot \tfrac23 \cdot \tfrac45 = \tfrac{4}{15}$ of whole numbers avoid all three primes, and $1000 \cdot 4/15 \approx 266.7$. The same pattern repeats every 30 numbers, with 8 survivors in each block.

\iqwrong{Multiplying $500 \cdot \tfrac23 \cdot \tfrac45$ with rounding at each step. It happens to be close, but a rounded product is not a count: inclusion-exclusion gives the exact number.}
\iqfollow{How would you code it for many primes?}{Sieve up to $n$, or walk the $2^k$ subsets of $k$ primes with the sign of each subset. The subset walk is the better choice when $n$ is huge and $k$ is small.}
\end{solution}

\begin{solution}{iq:iv1:counting-and-symmetry:11}
The probability is $150/243 = 50/81 \approx 0.617$.

Use inclusion-exclusion\index{inclusion-exclusion} over the empty boxes. There are $3^5 = 243$ assignments. Those missing a particular box number $2^5$, and there are three boxes. Those missing two particular boxes number just 1, and there are three pairs. So the count of assignments with no empty box is
\[
3^5 - 3 \cdot 2^5 + 3 \cdot 1 = 150.
\]
As a check, count by shapes. The box sizes are 3, 1, 1 or 2, 2, 1. There are $3 \cdot \binom{5}{3} \cdot 2 = 60$ of the first shape and $3 \cdot \binom{5}{2}\binom{3}{2} = 90$ of the second, and $60 + 90 = 150$.

\iqwrong{Starting by placing one ball in each box and letting the rest go anywhere: $5 \cdot 4 \cdot 3 \cdot 3^2 = 540$, which is more than the total. Any ball could have been the ``first'' in its box, so each assignment is counted several times.}
\iqfollow{What else does this compute?}{It is the same as asking for the chance that a coupon collector with 3 coupons finishes within 5 draws.}
\end{solution}

\begin{solution}{iq:iv1:counting-and-symmetry:12}
No letter correct has probability $44/120 = 11/30 \approx 0.367$, exactly two correct has probability $20/120 = 1/6$, and the first approaches $1/e \approx 0.368$.

A permutation with no fixed point is a derangement\index{derangement}. By inclusion-exclusion\index{inclusion-exclusion},
\[
D_n = n! \sum_{k=0}^{n} \frac{(-1)^k}{k!}, \quad\text{so}\quad D_5 = 44, \quad \frac{D_n}{n!} \to e^{-1}.
\]
The convergence is very fast. Even five letters are within $0.002$ of the limit.

For exactly two correct, choose which two, and derange the other three: $\binom{5}{2} D_3 = 10 \cdot 2 = 20$.

\iqwrong{Thinking the chance of no match must grow or shrink strongly with more letters. There are more letters to match, but each has less chance of matching, and the two effects cancel: the expected number of matches is always one, and the number of matches tends to a Poisson with mean 1.}
\iqfollow{Is exactly four correct possible?}{No. If four letters are right, the fifth envelope is the only one left, and the letter must be right too.}
\end{solution}

\enlargethispage{-3\baselineskip}%
\begin{solution}{iq:iv1:counting-and-symmetry:13}
The probability is $(6 - 4)/(6 + 4) = 1/5$.

This is the ballot theorem\index{ballot theorem}, and the proof uses reflection\index{reflection principle}. Of the $\binom{10}{4} = 210$ orders, count the bad ones, where the count is tied at some point. An order that starts with a B vote is bad for sure, since A has to catch up, which leaves a tie. There are $\binom{9}{3} = 84$ such orders.

Now take an order that starts with A and has a tie. Swap A and B votes up to the first tie. This gives an order starting with B, and the swap can be undone. So bad orders starting with A also number 84. The good orders number $210 - 168 = 42$, and $42/210 = 1/5$.

\iqwrong{Treating it as ``A has more votes, so is usually ahead'' and giving a large number. A lead from the very first vote is a strong demand, and it fails whenever B leads early.}
\iqfollow{And with 60 votes against 40?}{The chance is the same formula: $20/100 = 1/5$. Only the margin as a share of the votes matters, which tells you that the size of the count does not.}
\end{solution}

\begin{solution}{iq:iv1:counting-and-symmetry:14}
With five of each, the probability is $1/6$. With six \$5 and four \$10 notes, it is $3/7$.

Write a \$5 note as an up step and a \$10 note as a down step. The till has enough change exactly when the walk never goes below 0. Use reflection\index{reflection principle} on the bad orders: those that touch $-1$. Reflect the part of the walk up to the first touch of $-1$. It then starts from $-2$ and ends at the same point, so bad orders match every path from $-2$ to the end point. For five and five, such a path has six up steps and four down steps, so there are $\binom{10}{4} = 210$ of them, against the total of $\binom{10}{5} = 252$. So 42 are good, which is the Catalan number\index{Catalan numbers} $C_5$, and $42/252 = 1/6$.

The general result for $a$ fives and $b$ tens is $(a + 1 - b)/(a + 1)$. For six and four, $3/7$: $90$ of the 210 orders.

\iqwrong{Assuming that with equal notes the chance is about one half. Having enough change at the very end is automatic, and it is the path before the end that fails.}
\iqwhy{Reflection is the tool behind pricing a barrier option in a binomial model: a price path that crosses a level is matched with a reflected one that ends elsewhere.}
\end{solution}

\begin{solution}{iq:iv1:counting-and-symmetry:15}
The probability is exactly $1/2$.

Compare after the first ten flips of each. Let $X$ be your heads in the first 10 and $Y$ mine. By symmetry\index{symmetry argument}, $P(X > Y) = P(X < Y) = a$, say, and $P(X = Y) = 1 - 2a$. You win if $X > Y$, or if $X = Y$ and your last coin is heads. If $X < Y$, even a final head gives at most a tie, and a tie is not strictly more. So your chance is $a + \tfrac12(1 - 2a) = \tfrac12$.

A neater version: you have either more heads than I have, or more tails, and never both, because your 11 flips exceed my 10. For a fair coin, heads and tails have the same role, so each has probability $1/2$.

\iqwrong{Thinking the extra coin gives you more than half. It does give you the tie-break, but ties were counted against you to begin with, so the gain only just brings you to one half.}
\iqfollow{And with a biased coin?}{The symmetry fails. If heads is more likely than tails, then you win more than half the time, because the event ``more heads'' is favoured.}
\end{solution}

\begin{solution}{iq:iv1:counting-and-symmetry:16}
They are equally likely: each has probability exactly $1/52$.

Build the shuffled deck in two steps. Take out the card you care about, and shuffle the other 51. Now insert it at a random place, one of the 52 possible gaps\index{insertion argument}.

For the two of clubs: it is ``right after the first ace'' for exactly one of the 52 gaps, whichever the other 51 cards look like. So its chance is $1/52$.

For the ace of spades: use the same deal. The other 51 cards contain three aces, and the first of them is somewhere. If you insert the ace of spades right after it, it is the next card after the first ace. If you insert it earlier, it becomes the first ace itself, and it does not follow one. If you insert it later, it does not come next. So exactly one gap works, and again the chance is $1/52$.

\iqwrong{Saying the ace of spades is less likely, because it might have already gone as the first ace, or more likely, because aces cluster. Both feelings are correct about something, but the insertion count settles it.}
\iqwhy{Building a random object one card at a time from a simpler one often avoids a heavy conditioning, and interviewers look for this kind of step.}
\end{solution}

\begin{solution}{iq:iv1:counting-and-symmetry:17}
The probability is $\frac{39 \cdot 26 \cdot 13}{51 \cdot 50 \cdot 49} \approx 0.105$.

Place the aces one by one, thinking of the deck as 52 slots in four hands\index{sequential conditioning}. The first ace goes anywhere. The second must avoid the 12 other slots in its hand, so it has 39 good slots of the 51 left. The third has 26 good slots of the 50 left, and the fourth 13 of the 49 left.

The other route is by counting whole deals. The number of deals where each hand has an ace is $4!\binom{48}{12}\binom{36}{12}\binom{24}{12}$, against $\binom{52}{13}\binom{39}{13}\binom{26}{13}$ deals in all. It reduces to the same fraction, with a good deal more cancelling.

\iqwrong{The independent answer, $4!/4^4 \approx 0.094$. This treats each ace as choosing a hand independently, and it misses the fact that a hand holding an ace has less room for the next.}
\iqfollow{Which method carries further?}{Placing special cards one by one works whenever the special cards are few. Counting deals is better when you must control many cards at once.}
\end{solution}

\begin{solution}{iq:iv1:counting-and-symmetry:18}
Both sums have 27 outcomes, and sum 14 on four dice has 146.

Listing unordered triples for 10 gives $(1,3,6)$, $(1,4,5)$, $(2,2,6)$, $(2,3,5)$, $(2,4,4)$, $(3,3,4)$, with $6, 6, 3, 6, 3, 3$ orders, which is 27. They are equal by a symmetry\index{symmetry argument} of the die: replacing each face $x$ by $7 - x$ is a bijection of outcomes that sends a sum $s$ to $21 - s$, and $21 - 10 = 11$.

For four dice, use stars and bars\index{stars and bars} with inclusion-exclusion\index{inclusion-exclusion}. Writing $d_i = 1 + y_i$, you need $y_1 + y_2 + y_3 + y_4 = 10$ with $y_i \le 5$. With no ceiling that has $\binom{13}{3} = 286$ solutions. If one $y_i \ge 6$, take away 6 from it and count solutions of a sum to 4: $\binom{7}{3} = 35$, for each of 4 dice. Two cannot exceed at once, because $12 > 10$. So $286 - 140 = 146$.

\iqwhy{Stars and bars with a ceiling is the same calculation as a discrete convolution, and the mirror symmetry halves the work whenever you compute a dice table by hand.}
\end{solution}

\begin{solution}{iq:iv1:counting-and-symmetry:19}
There are $5775$ splits, and the chance that the two end up together is $3/11$.

For the number of splits, start with $12!$ orderings, cut into three blocks of four, and divide out the orders inside each block and the orders of the blocks: $12!/(4!^3\,3!) = 5775$. Alternatively, choose the 3 companions of a fixed person, then the 3 companions of the first person left over: $\binom{11}{3}\binom{7}{3} = 165 \cdot 35 = 5775$.

The second part needs no counting at all. Stand with the first person. Their group has 3 other places, and by symmetry\index{symmetry argument} the second person is equally likely to be in any of the other 11 places. So the chance is $3/11$. Counting confirms it: $\binom{10}{2}\binom{7}{3} = 1575$ of the 5775 splits have both together.

\iqwrong{Giving $1/3$ because there are three groups. That would be right if each person chose a group independently, but your own group has only three places left, while the others have four each.}
\iqfollow{With labelled groups?}{The number of splits becomes $12!/4!^3 = 34\,650$, six times as many. The probability stays the same.}
\end{solution}

\begin{solution}{iq:iv1:counting-and-symmetry:20}
The two players meet with probability $1/32$. The second-best player reaches the final with probability $32/63$.

For the first part, count matches rather than rounds. Exactly one player leaves in each match, so a 64-player tournament has 63 matches. There are $\binom{64}{2} = 2016$ pairs of players. By symmetry\index{symmetry argument} every pair is equally likely to be the pair in any given match, so each pair meets with probability $63/2016 = 1/32$, and in general $2/n$.

For the second part, the second-best can lose only to the best. The two are in opposite halves of the draw exactly when the second-best reaches the final. Given the place of the best player, there are 63 other places, of which 32 lie in the other half. So the chance is $32/63$.

\iqwrong{Working through the rounds, with a factor for each round that both must survive. That can be done and takes much longer. Linearity with the matches gives the answer in one line.}

\iqwhy{This is why the runner-up in a knockout is a poor estimate of the second-best: it is right just over half the time.}
\end{solution}

\begin{solution}{iq:iv1:counting-and-symmetry:21}
Six points give 31 regions, not 32, and 10 points give 256.

The formula is $1 + \binom{n}{2} + \binom{n}{4}$\index{regions of a disc}. Start with the disc, one region. Draw the chords one at a time. A new chord adds one region for itself, and one more for each chord it crosses, because each crossing cuts one more existing region in two. So the number of regions is $1 + (\text{chords}) + (\text{crossings})$. There are $\binom{n}{2}$ chords, and each set of 4 points gives exactly one crossing, the crossing of its two diagonals. So there are $\binom{n}{4}$ crossings. For $n = 6$: $1 + 15 + 15 = 31$. For $n = 10$: $1 + 45 + 210 = 256$.

The formula gives $1, 2, 4, 8, 16$ for $n = 1, \dots, 5$ only because $\sum_k \binom{n-1}{k}$ matches it for small $n$. A polynomial of degree 4 cannot keep up with $2^{n-1}$.

\iqwrong{Thirty-two, from the pattern. Interviewers ask this question to test whether a candidate will extrapolate from a few cases without a reason.}
\iqfollow{How can you check 31 with no formula?}{By Euler's formula. With $n + \binom{n}{4}$ vertices and $n + \binom{n}{2} + 2\binom{n}{4}$ edges, the count of inner faces is $E - V + 1$, which is the same result.}
\end{solution}

\begin{solution}{iq:iv1:counting-and-symmetry:22}
There are 292 ways.

Use dynamic programming\index{coin-change counting} on the coins, one denomination at a time. Let $w(a)$ be the ways to make amount $a$ with the coins used so far, starting from $w(0) = 1$. For each coin $c$, in turn, run $a$ upwards from $c$ and set $w(a) \leftarrow w(a) + w(a - c)$. Going upwards lets a coin be reused, and going coin by coin counts each combination once only, regardless of order. This is the same as reading a coefficient in the generating function $\prod_c 1/(1 - x^c)$.

By hand, work over the numbers of large coins. With no 50 and no 25, for instance, choosing $d$ dimes leaves $21 - 2d$ choices of nickels, and pennies fill the rest, which gives 121 ways for that case. Adding all the cases gives 292. An exhaustive loop over 50s, 25s, 10s and 5s, with pennies making up the rest, gives the same number.

\iqwrong{Looping amounts first and coins inside. That counts $1 + 5$ and $5 + 1$ as different, which counts orderings instead of combinations.}
\iqfollow{And the number of ways counting order?}{Swap the two loops. It grows enormous, because most of the sequences use many pennies.}
\end{solution}

\begin{solution}{iq:iv1:counting-and-symmetry:23}
There are 136 solutions.

With no ceiling, stars and bars\index{stars and bars} gives $\binom{32}{2} = 496$. Now subtract the solutions in which a variable exceeds 15. If $x \ge 16$, put $x' = x - 16$. Then $x' + y + z = 14$, which has $\binom{16}{2} = 120$ solutions. There are three such cases, and two variables cannot both be at least 16, since their sum would exceed 30. So inclusion-exclusion\index{inclusion-exclusion} stops there: $496 - 360 = 136$.

A check by geometry: the solutions form a triangular slice of a cube, and with a ceiling of half the total, the slice is the central triangle formed by the three midpoints. It has about a quarter of the unrestricted triangle's area, and $136/496 \approx 0.27$, a little more because of the boundary points.

\iqwhy{Ceilings appear whenever capacity limits a split. This is the standard way to treat them, and it stops after one or two terms when the ceiling is large.}
\end{solution}

\begin{solution}{iq:iv1:counting-and-symmetry:24}
The first is $1/70$, and the second is $1/34\,650$.

Ignore all the other cards. Only the relative order of the 8 aces and kings matters, and every way of placing those is equally likely\index{symmetry argument}. Of the $\binom{8}{4} = 70$ patterns of A and K, exactly one has all the A's first. The other 44 cards are irrelevant.

For twelve cards with the order aces, kings, queens, look at the patterns of three letters: there are $12!/(4!)^3 = 34\,650$ of them, and one is right.

\iqwrong{Multiplying conditional chances for each ace, such as $\tfrac{4}{52}$ for the first, then trying to track the deck. Removing irrelevant cards is the whole point.}

\iqwhy{This kind of reduction, where only the relative order of a few events matters, applies to any question of which event happens first among exchangeable times, such as which order fills first in a random arrival sequence.}
\end{solution}

\begin{solution}{iq:iv1:counting-and-symmetry:25}
The probability is $\binom{20}{10}/4^{10} = 184\,756/1\,048\,576 \approx 0.176$.

The direct sum is $\sum_k \binom{10}{k}^2 / 4^{10}$. The neat way sees the equal-heads event as a statement about 20 flips. Change the meaning of my flips by swapping heads with tails, which leaves them fair. Then ``your heads equals my heads'' becomes ``your heads plus my new heads equals 10'', because my new heads is $10$ minus my old heads. This is the event of exactly 10 heads in 20 fair flips, and its probability is $\binom{20}{10}/2^{20}$.

The same argument proves the identity\index{Vandermonde's identity} $\sum_k \binom{n}{k}^2 = \binom{2n}{n}$, and for large $n$ the probability is about $1/\sqrt{\pi n}$, here $0.178$.

\iqwrong{One over 11, since there are 11 possible counts. The counts are not equally likely, and the central ones carry most of the mass.}
\iqwhy{The chance of an exact tie after $n$ rounds falls only with the square root of $n$. Ties in a large sample are far more common than intuition suggests.}
\end{solution}

\begin{solution}{iq:iv1:counting-and-symmetry:26}
Containment has probability $(3/4)^{10} \approx 0.056$, disjointness has the same probability, and equality has probability $(1/2)^{10} = 1/1024$.

Each element acts independently\index{element-by-element counting}, and each has four cases: in $A$ only, in $B$ only, in both, or in neither. For $A \subseteq B$, only ``in $A$ only'' is forbidden, so each element contributes $3/4$. So there are $3^{10}$ good pairs among $4^{10}$.

For disjointness, only ``in both'' is forbidden, and again each element contributes $3/4$. Complementing $B$ turns one into the other, which is why they are equal. For equality, an element must be either in both or in neither, which has probability $1/2$.

\iqwrong{Conditioning on the size of $A$, and summing over that size. It does give the same answer, by the binomial theorem, but the element-by-element view avoids any sum.}
\iqfollow{What about $A \subseteq B \subseteq C$?}{Each element has 8 cases, of which 4 are allowed: $(1/2)^{10}$.}
\end{solution}

\enlargethispage{-3\baselineskip}%
\begin{solution}{iq:iv1:counting-and-symmetry:27}
In a row there are 816 choices, and in a circle there are 800.

In a row, use a gap argument\index{bijection}. Remove one chair from between each neighbouring pair of chosen chairs. Choosing 3 of 20 with no two adjacent then matches choosing 3 of $20 - 2 = 18$ chairs, with no condition. So $\binom{18}{3} = 816$.

In a circle, split on whether chair 1 is chosen. If it is, its two neighbours are out, and you need 2 non-adjacent chairs from the other 17 in a row: $\binom{16}{2} = 120$. If it is not, you need 3 non-adjacent from 19 in a row: $\binom{17}{3} = 680$. The total is 800. The general formula is $\frac{n}{n-k}\binom{n-k}{k}$.

The row and circle differ by those selections containing both end chairs, which are adjacent only on the circle. Fixing both, the third chair must come from chairs 3 to 18, which gives 16, and $816 - 16 = 800$ exactly.

\iqwhy{The difference between a row and a ring is a recurring trap. A clean way to check a circular answer is to count what closing the row into a ring removes.}
\end{solution}

\begin{solution}{iq:iv1:counting-and-symmetry:28}
The count is $1\,263\,360$ out of $10! = 3\,628\,800$, or about $0.348$.

Use inclusion-exclusion\index{inclusion-exclusion} over the set of couples that sit together. Fix $k$ couples and glue each into a block. There are then $10 - k$ objects, which can be arranged in $(10 - k)!$ ways, and each block has 2 internal orders, so the term is $\binom{5}{k} 2^k (10 - k)!$. So
\[
\sum_{k=0}^{5} (-1)^k \binom{5}{k} 2^k (10 - k)! = 3\,628\,800 - 3\,628\,800 + 1\,612\,800 - 403\,200 + 57\,600 - 3840.
\]
The first two terms cancel exactly, which is a small surprise, and the sum is $1\,263\,360$. The same formula gives 8 and 240 for two and three couples, which you can check by listing.

An approximation is a useful check. The expected number of couples together is $5 \cdot 2/10 = 1$, and if that number were Poisson\index{Poisson approximation} with mean 1, the answer would be $e^{-1} \approx 0.368$. The right answer is a bit lower, because the events are not independent, and it approaches $1/e$ for many couples.

\iqwrong{Multiplying $(1 - 1/5)^5 \approx 0.328$ as if the events were independent. Not bad as an estimate, but it is not exact, and an exact count was asked for.}
\iqfollow{How does a round table change it?}{Glued blocks arranged on a circle give $(9 - k)!$ arrangements with the same signs, and the count differs. Remember the factor for rotations.}
\end{solution}

\begin{solution}{iq:iv1:counting-and-symmetry:29}
There are 208 such subsets.

Use the roots-of-unity filter\index{roots-of-unity filter}. Let $\omega = e^{2\pi i/5}$, and consider $f(x) = \prod_{k=1}^{10}(1 + x^k)$. The coefficient of $x^s$ counts subsets with sum $s$. Averaging $f$ over the five fifth roots of unity keeps exactly the terms where $s$ is divisible by 5:
\[
N = \frac15 \sum_{j=0}^{4} f(\omega^j).
\]
At $j = 0$ you have $f(1) = 2^{10} = 1024$. For $j \ne 0$, the powers $\omega^{jk}$ run twice through all fifth roots of unity as $k$ goes from 1 to 10, so $f(\omega^j) = \bigl[\prod_{r=0}^{4}(1 + \omega^r)\bigr]^2$. Since $\prod_r (x - \omega^r) = x^5 - 1$, putting $x = -1$ gives $\prod_r (1 + \omega^r) = 2$. So each of the four terms is $2^2 = 4$, and
\[
N = \frac{1024 + 4 \cdot 4}{5} = 208.
\]
The others share the remainder evenly, at 204 for each non-zero residue. So the distribution is almost flat, and the small excess for residue 0 is exactly what the filter sees.

\iqwrong{Answering $1024/5$, which is not even a whole number, or assuming perfect uniformity. Each block of $\{1, \dots, 5\}$ holds a multiple of 5, so the residues are not perfectly balanced.}
\iqfollow{And for subsets of $\{1, \dots, 5\}$?}{The powers run once through the roots, so $(32 + 4 \cdot 2)/5 = 8$. Listing them confirms it.}
\end{solution}

\begin{solution}{iq:iv1:counting-and-symmetry:30}
There are 4 rings under rotation, and 3 if flipping is also allowed.

Use Burnside's lemma\index{Burnside's lemma}: the number of distinct rings is the average, over the symmetries, of the number of bead patterns that each symmetry leaves unchanged. There are $\binom{6}{3} = 20$ patterns.

Under rotation, the identity keeps all 20. A rotation by one or five steps moves all beads in a single cycle, so it would need all beads the same: 0 patterns. A rotation by two or four has two cycles of length 3, so one must be red and the other green: 2 patterns each. A rotation by three has three cycles of length 2, and red would have to be even: 0. The average is $(20 + 2 + 2)/6 = 4$. They are RRRGGG, RRGRGG, RRGGRG and RGRGRG.

Adding flips, three reflections through opposite beads each keep 4 patterns, because the two beads on the axis must differ and one of the other two pairs is red. The three through gaps keep none, because pairs make red even. The average is $(24 + 12)/12 = 3$. The two mirror images RRGRGG and RRGGRG become one.

\iqwrong{Dividing $20$ by the 6 rotations, which gives $10/3$. A pattern with symmetry has a smaller orbit, so simple division is wrong, and the lemma exists to correct it.}
\end{solution}

\begin{solution}{iq:iv1:counting-and-symmetry:31}
There are 10 cubes with two colours, and 57 with three colours.

Use Burnside's lemma\index{Burnside's lemma} on the 24 rotations of the cube. A rotation fixes a colouring exactly when each cycle of faces it creates is one colour, so with $k$ colours a rotation with $c$ cycles fixes $k^c$ colourings. The rotations split by cycle count:
\begin{itemize}
\item the identity has 6 cycles (one rotation);
\item a half-turn about the axis through two opposite faces has 4 cycles (three rotations);
\item a quarter-turn about that axis has 3 cycles (six rotations);
\item a half-turn about the axis through two opposite edges has 3 cycles (six rotations);
\item a rotation about the axis through two opposite corners has 2 cycles (eight rotations).
\end{itemize}
This totals $1 + 3 + 6 + 6 + 8 = 24$, and the number of cubes is
\[
\frac{k^6 + 3k^4 + 12k^3 + 8k^2}{24},
\]
which gives 10 for $k = 2$ and 57 for $k = 3$. With $k = 1$ it gives 1, a useful check on the group count.

For $k = 2$ you can list by number of faces of colour A, from 0 to 6: $1, 1, 2, 2, 2, 1, 1$. These add up to 10, confirming the formula.

\iqwrong{$2^6/24$, which is not even a whole number. Colourings with symmetry have smaller orbits.}
\iqfollow{And colourings using all six colours once each?}{No rotation except the identity fixes one, so $6!/24 = 30$.}
\end{solution}

\begin{solution}{iq:iv1:counting-and-symmetry:32}
A right angle has probability $60/220 = 3/11$. Strictly containing the centre has probability $40/220 = 2/11$.

There are $\binom{12}{3} = 220$ triangles. By the angle in a semicircle\index{angle in a semicircle}, a triangle with vertices on a circle has a right angle exactly when one side is a diameter. There are 6 diameters and 10 choices of third corner, which gives 60.

Go round the polygon and write the three arcs between consecutive corners as whole numbers adding up to 12. A triangle strictly contains the centre exactly when each arc is below 6. An arc of exactly 6 is the right-angled case, and an arc above 6 leaves the centre outside the triangle (an obtuse triangle).

Count the obtuse ones. Fix the corner at which the largest arc ends, so that the other two corners lie within the next 5 steps: there are 12 choices for that corner and $\binom{5}{2} = 10$ choices for the other two. So there are 120 obtuse triangles, and $220 - 60 - 120 = 40$ acute.

\iqwrong{Using $1/4$, which is the answer for three uniform points on a circle. A finite polygon has right-angled triangles that take up probability, so the acute share is lower. As the number of corners grows, the answer does approach $1/4$.}
\iqfollow{Three uniform points: why $1/4$?}{The arcs are a uniform point on a simplex, and the region where all arcs are below a half is the central triangle, a quarter of the simplex.}
\end{solution}

\begin{solution}{iq:iv1:counting-and-symmetry:33}
Use the pigeonhole principle\index{pigeonhole principle}, comparing the number of subsets with the number of possible sums. For 7 numbers the count of subsets is too small to force a match, and the powers 1, 2, 4, 8, 16, 32, 64 have all their 127 subset sums different.

Any 10 numbers from 1 to 100 have $2^{10} - 1 = 1023$ non-empty subsets. Each has a sum between 1 and at most $91 + 92 + \dots + 100 = 955$. There are more subsets than sums, so two different subsets have the same sum.

Those two may overlap. Remove their common elements from both: the sums stay equal to each other, and they stay different sets. Neither is left empty, because the numbers are positive, and if one set were emptied, the other would have to have sum zero. So you have two disjoint non-empty groups with the same sum.

For 7 numbers there are only 127 subsets, and sums can reach up to 679, so the pigeonhole argument gives nothing. The powers of two are a counterexample, because binary digits are unique. So the statement itself is false at 7, not just the proof.

\iqwrong{Trying to construct the two groups. The argument proves they exist without showing which they are, and that is the point of the question.}
\iqwhy{Pigeonhole arguments give very quick bounds, and the same comparison of counts gives a lower limit on the number of bits needed by a hash.}
\end{solution}

\begin{solution}{iq:iv1:counting-and-symmetry:34}
The probability is $232/1024 = 29/128 \approx 0.227$.

Use the reflection principle\index{reflection principle}. A path that reaches $+4$ and ends below 4 can be flipped after its first touch of 4, and that gives a path ending above 4. This is a bijection. So
\[
P(\max \ge 4) = P(S_{10} \ge 4) + P(S_{10} > 4).
\]
Paths ending at or above 4 must have 7 or more up steps: $\binom{10}{7} + \binom{10}{8} + \binom{10}{9} + \binom{10}{10} = 120 + 45 + 10 + 1 = 176$. Paths ending above 4 must have 8 or more: 56. The total is $176 + 56 = 232$ out of $2^{10}$.

\iqwrong{Doubling $P(S_{10} \ge 4)$, to get $352/1024$. The continuous-time formula does give a factor of two, but here the paths that end exactly on 4 must be counted once, not twice. That makes a big difference over only 10 steps.}
\iqfollow{What does the continuous limit say?}{Brownian motion over time $T$ reaches a level $a$ with probability $2P(W_T > a)$, where exact hits have no mass. The difference falls with the number of steps.}
\iqwhy{This is the reflection used for barrier options: the chance of touching a level in a period is about twice the chance of ending beyond it. Near the level on a coarse grid, the discrete correction matters.}
\end{solution}

\enlargethispage{-3\baselineskip}%
\begin{solution}{iq:iv1:counting-and-symmetry:35}
There are $55\,252$ lucky tickets.

The trick is a bijection\index{bijection}. Replace each of the last three digits $d$ by $9 - d$. A ticket is lucky when $a + b + c = d + e + f$. After the replacement, $d' + e' + f' = 27 - (d + e + f)$, so the condition becomes $a + b + c + d' + e' + f' = 27$. So you are counting six-digit strings whose digits add up to 27, which is a single coefficient: that of $x^{27}$ in $(1 + x + \dots + x^9)^6$.

Find this using stars and bars\index{stars and bars} with inclusion-exclusion\index{inclusion-exclusion}. With no ceiling, the number of solutions is $\binom{32}{5} = 201\,376$. Subtract 6 times the number with one digit at least 10, which is $\binom{22}{5} = 26\,334$, and add back 15 times the number with two such digits, $\binom{12}{5} = 792$. Three digits at least 10 would need 30, so you stop there:
\[
201\,376 - 158\,004 + 11\,880 = 55\,252.
\]
Check by the direct sum of squares: for each value $s$ from 0 to 27, the number of three-digit strings with sum $s$ is squared and these are added up, which gives the same total. The peak term is $75^2$ at $s = 13$.

\iqwrong{Assuming that about 1 in 28 tickets is lucky, which is 35,714. Sums near the centre are much more likely, so a match is more likely than that: the true rate is about 1 in 18.}
\end{solution}

\begin{solution}{iq:iv1:counting-and-symmetry:36}
There are $2^{16} = 65\,536$ such tables.

Think of free choices\index{degrees of freedom}. Fill in the top-left $4 \times 4$ block in any way: that is $2^{16}$ choices. Each of the first four rows then has its last entry forced so that the row has even sum, and each of the first four columns has its last entry forced in the same way. Only the bottom-right corner remains.

The corner must fix both the last row and the last column, and these two demands agree. The sum of the last column's top four entries and the sum of the last row's first four entries are each the parity of the sum of the $4 \times 4$ block: each is obtained by adding up the rows, or the columns, of the block. So the corner exists and is unique, whichever block you chose.

In general the count for $n \times n$ is $2^{(n-1)^2}$. Enumeration gives 2, 16 and 512 for $n = 2, 3, 4$, as it should.

\iqwrong{Counting equations and unknowns naively: 25 unknowns less 10 constraints leaves $2^{15}$. The 10 equations are not independent, because the sum of all rows and the sum of all columns are the same sum, so one equation is redundant.}
\iqwhy{Parity grids are the idea behind RAID storage and simple error-correcting codes. The redundant equation is what lets you find a failed entry.}
\end{solution}

\enlargethispage{-3\baselineskip}%
\begin{solution}{iq:iv1:counting-and-symmetry:37}
The count is $2^n - n - 1$, which is 26 for $n = 5$.

A permutation with exactly one ascent is two descending runs placed one after the other\index{Eulerian numbers}. Each is decided by its set of entries, because a descending run has only one order. So choose a subset $S$ for the first run: there are $2^n$ subsets, and the other entries form the second run.

Each permutation with one ascent arises from exactly one subset. But some subsets do not make an ascent. If the first run's last entry is larger than the second run's first entry, the whole thing descends, and there is no ascent. That happens exactly when $S$ is the set of the top $k$ values, for $k = 0, 1, \dots, n$: that is $n + 1$ subsets, each giving the single fully descending permutation. So the count is $2^n - (n + 1)$.

Enumeration agrees: $n = 3$ gives 4, namely 132, 213, 231 and 312, and $n = 4$ gives 11, $n = 5$ gives 26, $n = 6$ gives 57.

\iqwrong{Choosing where the ascent goes, with $n - 1$ positions, and a subset of fixed size. That gives the right total only if you subtract the subsets at each size that make no ascent, and it is a slower route to the same count.}
\iqfollow{What are these numbers called?}{Eulerian numbers. They count permutations by their number of ascents, and they appear when you sum powers of integers.}
\end{solution}

\begin{solution}{iq:iv1:counting-and-symmetry:38}
The pair is $\{1, 2, 2, 3, 3, 4\}$ and $\{1, 3, 4, 5, 6, 8\}$, known as Sicherman dice\index{Sicherman dice}, and it is the only other pair.

Use generating functions\index{generating functions}. A die is a polynomial $D(x) = \sum x^{\text{face}}$, and the sum of two dice has the product as its generating function. The target is
\[
(x + x^2 + \dots + x^6)^2 = x^2 (1 + x)^2 (1 + x + x^2)^2 (1 - x + x^2)^2,
\]
and these factors are irreducible over the integers. Polynomial factorisation is unique, so each die $D$ is a product of some of these. Two conditions apply to each die. First, it must have $D(0) = 0$, because the faces are positive, so it takes one factor of $x$. Second, it must have $D(1) = 6$, because it has six faces. At $x = 1$, the factors take the values $2$, $3$ and $1$. To get 6, each die must take exactly one $(1 + x)$ and one $(1 + x + x^2)$.

Only the $(1 - x + x^2)$ factors, whose value at 1 is 1, can move. Sharing them $1 + 1$ gives standard dice. Sharing them $2 + 0$ gives
\[
x(1 + x)(1 + x + x^2) = x + 2x^2 + 2x^3 + x^4
\]
and $x(1 + x)(1 + x + x^2)(1 - x + x^2)^2 = x + x^3 + x^4 + x^5 + x^6 + x^8$.

A computer search of all face sets confirms it: there are exactly these two solutions.

\iqwrong{Searching by hand with ``the two dice must sum to 42 together''. It is a true necessary condition, but it does not lead you to the pair, and it proves nothing about uniqueness.}
\iqfollow{Three dice?}{The same factorisation has more possibilities, and there are non-standard triples, also found by moving the $(1 - x + x^2)$ factors.}
\end{solution}

\begin{solution}{iq:iv1:counting-and-symmetry:39}
There are 312 seatings, and the chance that a random alternating seating qualifies is $13/120 \approx 0.108$.

This is the m\'enage problem\index{m\'enage problem}. Fix one man to remove the rotations: the other men can be placed in $4! = 24$ ways. Between men there are five places for the women, and each woman is barred from two adjacent places: the ones beside the matching partner. So you need the number of assignments of 5 women to 5 places, with these barred positions, and this number does not depend on how the men sit.

The bars form a single ring of 10 cells: woman, place, woman, place, and so on. Choosing $k$ violations is choosing $k$ cells of the ring with no two neighbouring, because a woman or a place cannot be used twice. There are $\frac{2n}{2n - k}\binom{2n - k}{k}$ such choices, and the $n - k$ others are placed freely, so by inclusion-exclusion\index{inclusion-exclusion},
\[
M_n = \sum_{k=0}^{n} (-1)^k \frac{2n}{2n - k}\binom{2n - k}{k} (n - k)!, \qquad M_5 = 13.
\]
The total is $4! \cdot 13 = 312$, out of $4! \cdot 5! = 2880$ alternating seatings. Checking all 2,880 by machine gives the same count, and the formula also gives $M_6 = 80$.

\iqwrong{Treating each woman as independently having a $3/5$ chance of a valid place, which gives $(3/5)^5 \approx 0.078$. The choices compete for places, and they are far from independent.}
\iqfollow{For large $n$?}{Each woman has 2 barred places out of $n$, so the expected number of violations is 2, and the fraction tends to $e^{-2} \approx 0.135$, which is a fair check here.}
\end{solution}

\begin{solution}{iq:iv1:counting-and-symmetry:40}
Each prisoner should start at the box with their own number and follow the numbers inside. The chance of success is $1 - \sum_{k=51}^{100} 1/k \approx 0.312$, compared with $2^{-100}$ for random guessing.

Each box holds a number that is the label of another box, so the boxes form a permutation\index{cycle structure of permutations}. If a prisoner opens box $i$, finds the number $j$, then opens box $j$ and so on, the prisoner walks along the cycle that contains $i$. That cycle ends at the box holding $i$, which is exactly what the prisoner wants. So the prisoner succeeds if and only if the cycle is at most 50 long. All succeed if and only if the permutation has no cycle longer than 50.

A permutation of 100 has at most one cycle longer than 50. The number of permutations with a cycle of length exactly $k > 50$ is found by choosing its elements, $\binom{100}{k}$, arranging them in a cycle, $(k - 1)!$, and arranging the rest, $(100 - k)!$. This is $100!/k$. So the chance of a long cycle is $\sum_{k=51}^{100} 1/k \approx \ln 2$, and success has probability about $1 - \ln 2 \approx 0.307$.

The strategy works because success is highly correlated: all prisoners succeed or fail together, depending on one property of the permutation.

\iqwrong{``Each has $1/2$, so it is $2^{-100}$ whatever they do, because nobody can pass information.'' The individual chance really is $1/2$. What the plan changes is the correlation between them, and independence is what is being given up.}
\iqwhy{It is the best example of a pattern that matters for risk: you cannot change a position's marginal behaviour, but correlation decides what happens to all the positions together.}
\end{solution}
